\chapter[Materials, Methods and Implementation]{Materials, Methods and \\ Implementation}
\label{chap:methods}

Translating the theoretical framework of Chapter~\ref{chap:introduction} into a feasible empirical analysis requires addressing three practical challenges inherent to clinical stereo-electroencephalography:
\begin{itemize}
    \item the computational overhead of handling multi-channel recordings of several gigabytes;
    \item the contamination of intracranial potentials by extracranial polygraphic artifacts;
    \item the non-stationarity and inter-patient variability of clinical metadata.
\end{itemize}
This chapter details the dataset structure, the signal preprocessing pipeline, and the core numerical algorithms implemented in Python to extract and calibrate the dynamic functional networks.

\section{The sEEG Dataset and Memory-Mapped Extraction}
\label{sec:dataset_extraction}

The empirical dataset analyzed in this study consists of invasive multi-channel sEEG recordings acquired from patients with drug-resistant focal epilepsy. Each recording is stored as a compressed NumPy archive (\texttt{.npz}) containing the continuous voltage matrix $\mathbf{X} \in \mathbb{R}^{N_{\text{raw}} \times M_{\text{tot}}}$, sampled at $f_s = 1000\text{ Hz}$, alongside an array of clinical channel labels and a metadata table specifying the annotated seizure onset time $t_{\text{onset}}$.

In the raw archives, the total number of recorded traces ranges between $N_{\text{raw}} \approx 170$ and $250$ per patient. However, this raw count includes numerous auxiliary and polygraphic reference channels. The actual number of implanted intracranial contacts, which define the nodes of our functional network, ranges between $N \approx 110$ and $170$ valid channels per patient (typically denoted by alphabetical prefixes such as \texttt{A}, \texttt{B}, \texttt{H}, \texttt{S'}, \texttt{Y'}).

\subsection{Out-of-Core Slicing via Memory Mapping}
Because continuous clinical monitoring spans hours of activity at kilohertz resolution, loading an entire raw recording array directly into volatile memory (RAM) is computationally inefficient and can lead to memory overflow. However, the critical dynamical reorganization governing the transition from the inter-ictal baseline to the ictal state unfolds over a temporal horizon of a few minutes around seizure onset.

To isolate the phase transition with minimal memory footprint, we access the raw archives in read-only memory-mapped mode (\texttt{mmap\_mode="r"}). This operating system mechanism maps the byte offsets of the uncompressed array directly onto virtual address space without transferring the full file into RAM. Given the aligned seizure onset timestamp $t_{\text{onset}}$, we slice a symmetric $240\text{ s}$ transition window spanning $\Delta t_{\text{pre}} = 120\text{ s}$ prior to onset and $\Delta t_{\text{post}} = 120\text{ s}$ following onset, shifting the temporal axis so that $t = 0\text{ s}$ corresponds exactly to the transition point.

\subsection{Artifact Channel Filtering and Channel-Wise Standardization}
Clinical sEEG systems routinely record auxiliary polygraphic channels alongside intracranial depth electrodes, labeled with prefixes such as \texttt{POL} (polygraphy/polarity references), \texttt{ECG} (electrocardiogram), or \texttt{MK} (event markers). During the pre-ictal and ictal phases, autonomic changes and even subtle patient motor activity can generate high-amplitude electrical fluctuations on these extracranial channels.

If retained in the graph, such signals may exhibit strong (anti-)correlations with cortical activity, leading to false positives in network analysis. We therefore filter out all non-intracranial labels prior to network construction. Subsequently, because electrode impedance and local tissue conductivity vary across recording sites, raw voltage amplitudes cannot be compared directly without normalization.

Each retained intracranial channel $i \in \{1, \dots, N\}$ is independently standardized over the extracted $240\text{ s}$ chunk via a channel-wise $z$-score transformation:
\begin{equation}
    z_i(t) = \frac{x_i(t) - \mu_i}{\sigma_i}
    \label{eq:zscore}
\end{equation}
where $\mu_i$ and $\sigma_i$ denote the temporal mean and standard deviation of channel $i$ across the extracted transition window. The core Python implementation is shown below:

\begin{lstlisting}[language=Python, caption={Memory-mapped transition window extraction, extracranial channel filtering, and channel-wise $z$-score normalization.}]
def extract_transition_chunk(file_path, onset_sec, pre_sec=120.0, post_sec=120.0):
    # Access large archive on disk without loading full array into RAM
    data = np.load(file_path, mmap_mode="r")
    sfreq = float(data["sfreq"])
    signal_arr = data[data.files[0]]
    channel_names = data["channels"]

    if signal_arr.shape[0] > signal_arr.shape[1]:
        signal_arr = signal_arr.T  # Enforce (N_channels, N_samples)

    start_idx = max(0, int((onset_sec - pre_sec) * sfreq))
    end_idx = min(signal_arr.shape[1], int((onset_sec + post_sec) * sfreq))

    # Pull only the target [-120s, +120s] window into memory
    raw_chunk = np.array(signal_arr[:, start_idx:end_idx], dtype=np.float64)

    # Filter out extracranial/polygraphic channels (POL, ECG, MK)
    if len(channel_names) > 0:
        valid_idx = [
            i for i, ch in enumerate(channel_names)
            if not any(tag in str(ch).upper() for tag in ("POL", "ECG", "MK"))
        ]
        raw_chunk = raw_chunk[valid_idx, :]
        channel_names = channel_names[valid_idx]

    # Channel-wise z-score standardization
    means = np.mean(raw_chunk, axis=1, keepdims=True)
    stds = np.std(raw_chunk, axis=1, keepdims=True)
    stds[stds == 0.0] = 1.0
    normalized_chunk = (raw_chunk - means) / stds

    rel_time = (np.arange(start_idx, end_idx) - int(onset_sec * sfreq)) / sfreq
    return normalized_chunk, rel_time, sfreq, channel_names
\end{lstlisting}

\section{Sliding-Window Tensor Construction}
\label{sec:tensor_construction}

To track the non-stationary topological evolution of the sEEG network across the $240\text{ s}$ window, the standardized signal matrix $\mathbf{Z} \in \mathbb{R}^{N \times M}$ is partitioned into overlapping sliding windows of fixed duration $T = 2.0\text{ s}$ ($2000\text{ samples}$) and step size $\Delta \tau = 0.5\text{ s}$ ($500\text{ samples}$). This parameterization yields a total of $N_\text{wins} = 477$ temporal frames centered at $\tau \in [-119.0\text{ s}, +119.0\text{ s}]$.

The choice of $T = 2.0\text{ s}$ represents an optimal physical trade-off between statistical stability and temporal resolution: a $2000$-sample window is sufficiently long to provide robust estimation of zero-lag Pearson correlations, while a $0.5\text{ s}$ stride resolves rapid pre-ictal decoupling and onset recruitment dynamics. Crucially, both $T$ and $\Delta \tau$ are held strictly constant across all patients to ensure that the variance of the sample correlation estimator remains invariant across the cohort.

Within each window $T_\tau$, we compute the signed Pearson correlation matrix $\mathbf{W}(\tau) \in [-1, 1]^{N \times N}$ according to Equation~(\ref{eq:pearson}), setting the diagonal elements to zero ($W_{ii}(\tau) = 0$) and clipping numerical floating-point inaccuracies to $[-1.0, 1.0]$. As established in Section~\ref{sec:theoretical_framework}, no absolute value transformation is applied, preserving the distinction between cooperative in-phase links ($W_{ij} > 0$) and competitive counter-phase links ($W_{ij} < 0$).

\begin{lstlisting}[language=Python, caption={Sliding-window computation of the signed dynamic functional connectivity tensor $\mathbf{W}(\tau)$.}]
def compute_dynamic_functional_connectivity(signal, sfreq, window_sec=2.0, step_sec=0.5):
    n_nodes, n_samples = signal.shape
    win_samples = int(window_sec * sfreq)
    step_samples = int(step_sec * sfreq)
    starts = np.arange(0, n_samples - win_samples + 1, step_samples)

    W = np.zeros((len(starts), n_nodes, n_nodes), dtype=np.float64)
    centers_sec = (starts + win_samples / 2.0) / sfreq

    for idx, s in enumerate(starts):
        segment = signal[:, s : s + win_samples]
        corr = np.corrcoef(segment)
        corr = np.nan_to_num(corr, nan=0.0)
        np.fill_diagonal(corr, 0.0)
        W[idx] = np.clip(corr, -1.0, 1.0)

    return W, centers_sec
\end{lstlisting}

\section{Data-Driven Calibration and Statistical Epochs}
\label{sec:calibration_epochs}

Applying network analysis across multiple clinical sEEG recordings requires addressing three methodological challenges widely documented in the brain network literature:
\begin{enumerate}
    \item[(i)] \textbf{Clinical vs.\ electrophysiological onset:} Hospital metadata frequently annotate the appearance of behavioral symptoms (\textit{clinical onset}) rather than the intracranial initiation of synchronous discharges (\textit{electrophysiological onset}) \cite{bartolomei2001dynamical, schindler2007assessing};
    \item[(ii)] \textbf{Implantation density and threshold bias:} Comparing graphs across subjects using a fixed correlation threshold introduces systematic bias when baseline connectivity density varies due to electrode implantation geometry \cite{vanwijk2010comparing, bullmore2009complex};
    \item[(iii)] \textbf{Inter-ictal spikes and frame-to-frame noise:} Subclinical inter-ictal epileptiform discharges (IEDs) during the pre-ictal period produce transient synchronization bursts that can corrupt baseline reference windows and single-frame estimates \cite{mormann2007seizure, jiruska2013synchronization}.
\end{enumerate}
Following established practices in functional network analysis, we address these factors through empirical standardization criteria prior to computing multiscale metrics.

\subsection{Electrophysiological Onset Alignment}
\label{subsec:onset_alignment}
In invasive sEEG monitoring, the earliest intracranial signature of a focal seizure often precedes clinical manifestations by several seconds \cite{bartolomei2001dynamical}. Because evaluating premonitory metrics in $[-15, 0]\text{ s}$ requires $t = 0\text{ s}$ to coincide with the exact point of intracranial symmetry breaking, relying on delayed clinical timestamps would contaminate the pre-ictal window with early seizure discharges.

Following energy-based and correlation-based onset detection approaches \cite{bartolomei2001dynamical, schindler2007assessing}, we align $t = 0\text{ s}$ by locating the sharpest upward step transition in both signal amplitude and network synchronization across a broad search window ($\pm 240\text{ s}$ around the nominal marker, accessed via \texttt{mmap\_mode="r"}). At each step $\tau$, we compute the root-mean-square voltage amplitude $E_{\text{RMS}}(\tau)$ across the top $15\%$ most active channels (capturing focal amplitude surges) and the reference cooperative link density $D^{+}_{0.50}(\tau)$ (fraction of pairs with $W_{ij}(\tau) \geq 0.50$). Both profiles are min-max normalized to $[0, 1]$ and averaged into a composite trace $\Phi(\tau) = \frac{1}{2}\hat{E}_{\text{RMS}}(\tau) + \frac{1}{2}\hat{D}^{+}_{0.50}(\tau)$. The electrophysiological onset $\tau^*$ is then identified via a moving-average step detector over a $\Delta H = 5.0\text{ s}$ ($10$-frame) horizon:
\begin{equation}
    \tau^* = \arg\max_{\tau} \left( \langle \Phi(t) \rangle_{t \in [\tau, \, \tau + \Delta H]} - \langle \Phi(t) \rangle_{t \in [\tau - \Delta H, \, \tau]} \right)
    \label{eq:onset_detector}
\end{equation}
Averaging over $5.0\text{ s}$ filters out brief inter-ictal spikes ($1\text{--}2\text{ s}$), whereas the sustained step transition into the seizure produces a clear global maximum $\tau^*$, defining the aligned origin $t = 0\text{ s}$.

\subsection{Fixed-Density Sub-Critical Threshold Calibration}
\label{subsec:threshold_calibration}
As shown by Van Wijk et al.\ \cite{vanwijk2010comparing} and Bullmore \& Sporns \cite{bullmore2009complex}, applying a single universal threshold $\theta$ across subjects with different spatial sampling densities confounds topological comparisons: densely spaced implants may already percolate ($S(\tau) \approx 1$) during resting state at $\theta = 0.50$, whereas sparser implants remain fragmented. To ensure cross-patient comparability, we adopt \textit{fixed edge-density calibration} \cite{vanwijk2010comparing}, selecting $\theta$ individually so that all baseline networks share a comparable sub-critical sparsity.

Specifically, we evaluate a discrete grid $\theta \in \{0.45, 0.50, 0.55, 0.60, 0.65, 0.70\}$ on each pre-extracted chunk and select the threshold whose mean cooperative link density during the quiet baseline matches $D^{+}_{\text{base}}(\theta) \approx 7.0\%$ (within the canonical $5\%\text{--}10\%$ sparse regime \cite{bullmore2009complex}). Analogously, for visualizing global negative density $D^{-}(\tau)$, we select $\theta_{\text{neg}} \in \{-0.30, -0.35, -0.40, -0.45\}$ to match a baseline anticorrelation density of $\sim 4.5\%$.

\subsection{Spike-Free Baseline Selection and Statistical Epochs}
\label{subsec:statistical_epochs}
To prevent subclinical inter-ictal bursts from biasing reference values \cite{mormann2007seizure, jiruska2013synchronization}, nodal observables are averaged over four temporal macro-epochs:
\begin{enumerate}
    \item \textbf{Quiet Baseline Epoch ($\mathcal{E}_{\text{base}}$, $50\text{ s}$ window within $[-115, -20]\text{ s}$):} Selected automatically by scanning contiguous $50\text{ s}$ ($100$-frame) intervals and minimizing an empirical variability score ($\mu_{D^+} + 2.5\,\sigma_{D^+} + 5.0\,\sigma_{\text{RMS}}$) to select intervals with as few isolated pre-ictal spikes as possible.
    \item \textbf{Premonitory Window ($\mathcal{E}_{\text{prem}} = [-15.0, 0.0]\text{ s}$):} Immediate pre-onset horizon over which counter-phase decoupling $\langle d_i^{(-)} \rangle_{\text{prem}}$ is averaged \cite{mormann2003epileptic, jiruska2013synchronization}.
    \item \textbf{Onset Recruitment Epoch ($\mathcal{E}_{\text{onset}} = [0.0, +15.0]\text{ s}$):} Early recruitment window over which eigenvector centrality $\langle v_{1,i} \rangle_{\text{onset}}$ is averaged to identify primary onset hubs.
    \item \textbf{Established Ictal Epoch ($\mathcal{E}_{\text{ictal}} = [+15.0, +60.0]\text{ s}$):} Mature seizure propagation phase.
\end{enumerate}
At a step of $\Delta \tau = 0.5\text{ s}$, each $15\text{ s}$ window ($\mathcal{E}_{\text{prem}}$ and $\mathcal{E}_{\text{onset}}$) averages $30$ consecutive frames, suppressing single-frame noise while remaining localized around $t = 0\text{ s}$.

\begin{lstlisting}[language=Python, caption={Step-gradient onset alignment and fixed-density baseline/threshold calibration.}, label={lst:calibration}]
def estimate_automatic_onset(abs_times, rms_energy, pos_sync, horizon_sec=5.0):
    # Min-max normalize energy and sync density to [0, 1]
    rms_norm = (rms_energy - np.min(rms_energy)) / (np.ptp(rms_energy) + 1e-8)
    sync_norm = (pos_sync - np.min(pos_sync)) / (np.ptp(pos_sync) + 1e-8)
    combined = 0.5 * rms_norm + 0.5 * sync_norm

    # Step difference between 5s post-window and 5s pre-window
    dt = abs_times[1] - abs_times[0] if len(abs_times) > 1 else 0.5
    h = max(1, int(horizon_sec / dt))
    jump_score = np.zeros_like(combined)
    for i in range(h, len(combined) - h):
        jump_score[i] = np.mean(combined[i : i + h]) - np.mean(combined[i - h : i])

    return abs_times[int(np.argmax(jump_score))], jump_score

def select_quiet_baseline_and_threshold(times, pos_densities, rms_energy):
    # 1. Scan [-115s, -20s] for the most stationary 50s (100-frame) baseline window
    win_len = int(50.0 / (times[1] - times[0]))
    pre_idx = np.where((times >= -115.0) & (times <= -20.0))[0]
    sync_ref = pos_densities[0.50]

    best_base, min_cost = (-90.0, -20.0), float("inf")
    for s in range(pre_idx[0], pre_idx[-1] - win_len + 1):
        seg_s, seg_r = sync_ref[s : s + win_len], rms_energy[s : s + win_len]
        cost = np.mean(seg_s) + 2.5 * np.std(seg_s) + 5.0 * np.std(seg_r)
        if cost < min_cost:
            min_cost = cost
            best_base = (round(float(times[s]), 0), round(float(times[s + win_len]), 0))

    # 2. Fixed-density calibration: select threshold closest to ~7.0% baseline sparsity
    mask = (times >= best_base[0]) & (times <= best_base[1])
    best_th = min(pos_densities.keys(), key=lambda th: abs(np.mean(pos_densities[th][mask]) - 7.0))
    return best_base, best_th
\end{lstlisting}